\chapter{Implementation}

% Number layout FUCKERY needed. Could really help!!

This chapter describes the implementation of the different components of the analysis tool, working from the naive 
baseline, exploring matrix representations and computations, continuing to data-driven optimisations, and concluding 
with acceleration of compute using additional hardware.

\section{Naive Baseline}

As described in Section \ref{intro:methodology}, classical techniques for 
ILP limit studies involve a sequential algorithm iterating through the program trace. We take inspiration from PerPI's 
use of the Intel Pin Tool \cite{perpi}. The algorithm presented successfully finds the longest path in the DDG described
in Section \ref{prep:register-only-chains}.

For each instruction in the program trace, the algorithm:
\begin{enumerate}
  \item  Examines static read dependencies, and computes the size of the largest chain of dependencies ending in a dependent register.
  \item  Examines static write dependencies, and propagates the longest chain of dependencies to the next iteration.
\end{enumerate}

\begin{lstlisting}[language=C++]
struct StaticDep {
  int[] reads;
  int[] writes;
}

function findLongestPath(int[] trace, map<int, StaticDep> deps) -> int {
  int[] curDeps = [0,...,0];

  int maxChainLength = 0;

  for (int instr : trace) {
    int chainLength = 1;

    /* Examine read dependencies */
    for (int reg : deps[instr].reads) {
      chainLength = max(chainLength, curDeps[reg].snd + 1);
    }

    maxChainLength = max(maxChainLength, chainLength);

    /* Propagate write dependencies */
    for (int reg : deps[instr].writes) {
      curDeps[reg].fst = true;
      curDeps[reg].snd = chainLength;
    }
  }

  return maxChainLength;
}
\end{lstlisting}

\subsection{Verifying the baseline}

%TODO: not sure i want to write like this.
We can verify a result from the algorithm by inspecting a small example.

The \verb|cksum| program trace (Section \ref{prep:datasets}) has a small inner loop of 
7 instructions, which iterates 65536 times each time it is invoked.

\begin{figure}[H]
  \centering
    \begin{minipage}{0.6\textwidth}
      \begin{lstlisting}[frame=none,numbers=none,escapeinside={(*}{*)}, xleftmargin=4em]
        (*\offset{...}*)
        (*\offset{1c68}*) (*\color{gray}\textit{eor \marker{x19_1}x19, x0, x19, lsl \#8}*)
        (*\offset{1c6c}*) (*\color{gray}\textit{cmp x2, x3}*)
        (*\offset{1c70}*) (*\color{gray}\textit{b.ne 1c58}*)
        (*\offset{1c58}*) ldrb (*\marker{w1_1}*)w1, [w2], #1
        (*\offset{1c5c}*) eor (*\marker{x1_1}*)x1, x1, (*\marker{x19_2}*)x19, lsr #24
        (*\offset{1c60}*) and (*\marker{x1_3}*)x1, (*\marker{x1_4}*)x1, #0xff
        (*\offset{1c64}*) ldr (*\marker{x0_1}*)x0, [x23, (*\marker{x1_5}*)x1, lsl #3]
        (*\offset{1c68}*) eor (*\marker{x19_3}*)x19, (*\marker{x0_2}*)x0, x19, lsl #8
        (*\offset{1c6c}*) cmp x2, x3
        (*\offset{1c70}*) b.ne 1c58
        (*\offset{1c58}*) (*\color{gray}\textit{ldrb w1, [w2], \#1}*)
        (*\offset{1c60}*) (*\color{gray}\textit{eor x1 , x1 , \marker{x19_4}x19 , lsr \#24}*)
        (*\offset{...}*)
      \end{lstlisting}
    \end{minipage}
    \begin{tikzpicture}[remember picture, overlay]
        % 1. x19 (eor) -> x19 (eor)
        \draw[deparrow] ($(x19_1)+(0.2,0)$) to ($(x19_2)+(0.05,0.1)$);

        % 2. x1 (eor) -> x1 (and)
        \draw[deparrow] ($(x1_1)+(0.2,0)$) to ($(x1_4)+(0.05,0.1)$);
        
        % 3. x1 (and) -> x1 (ldr)
        \draw[deparrow] ($(x1_3)+(0.35,0.05)$) to ($(x1_5)+(0.05,0.1)$);

        % 4. x0 (ldr) -> x0 (eor)
        \draw[deparrow] ($(x0_1)+(0.35,0.05)$) to ($(x0_2)+(0.05,0.1)$);

        % 5. x19 (eor) -> x19 (eor)
        \draw[gdeparrow] ($(x19_3)+(0.4,0)$) to ($(x19_4)+(0.05,0.1)$);
    \end{tikzpicture}
  \caption{Partially unrolled cksum loop body in Arm Assembly, annotated with critical path dependencies}
  \label{fig:cksum-deps}
\end{figure}


Figure \ref{fig:cksum-deps}
shows the Arm assembly code for cksum's inner loop, including instructions from previous and successive loop iterations.
Data dependencies from written registers to read registers are drawn on the figure in red, showing the critical path 
for one iteration of the loop. Note that the critical path repeats across loop iterations.

Given that there are 7 instructions in the loop body, and the length of the critical path over one body is 4, 
we should expect that the ILP limit is $7/4=1.75$. Executing the sequential algorithm over one invocation of the inner 
loop (65536 iterations) processes a total of 458752 instructions, and a critical path length of 262145. Considering this 
path length includes the outgoing edge from the last execution of instruction \verb|1c68|, we can subtract 1 from this.
Then, the ILP limit calculated is $458752/262144=1.75$, as expected.

\subsection{Parallelisation?}

%TODO: Rewrite - this is crappy explanation.
Since only the longest chain is stored at each iteration and all other chains discarded, the \verb|curDeps| array 
holding the longest chain length ending at each register is not a valid memoisable result for a small instruction window 
repeating many times in the algorithm. All the longest dependency chains starting at some register and
ending at another have to be encoded in the result, which pushes us towards the matrix representation.
Then, propagation of dependencies for all the chains is simply the multiplication of the matrices.

\section{Matrix Multiplications}

This section details several approaches to representing and multiplying matrices together on the CPU.

\subsection{A naive matrix implementation}

Matrices are usually represented by a 2-dimensional array of elements. Row-major order and column-major order are methods 
for storing multidimensional arrays in linear memory (arrays of a single dimension).

\begin{figure}[H]
  \centering
  \begin{tikzpicture}[
    % Custom styles for the grids and labels
    matgrid/.style={
        matrix of math nodes,
        nodes={anchor=center, minimum size=1.0cm, font=\Large},
        row sep=0.2cm,
        column sep=0.3cm,
    },
    title/.style={font=\bfseries, align=center},
    % Recreate the specific colors from the reference image
    rm_color/.style={red!60}, % Salmon-red for Row-Major
    cm_color/.style={cyan}  % Cyan for Column-Major
]

% ==========================================
% LEFT: ROW-MAJOR TRAVERSAL
% ==========================================
% 1. The Matrix Grid (inside the brackets)
\matrix (RM) [matgrid] {
    a_{00} & a_{01} & a_{02} \\
    a_{10} & a_{11} & a_{12} \\
    a_{20} & a_{21} & a_{22} \\
};

% 2. The Title (Placed relative to the matrix)
\node[title, above=0.3cm of RM-1-2] (rm_title) {Row-major order};

% 3. The Math Brackets (drawn manually for precision)
\draw[line width=1.5pt] ($(RM-1-1.north west)+(-0.2,-0.1)$) -- ++(-0.2,0) |- ($(RM-3-1.south west)+(-0.2,0.1)$);
\draw[line width=1.5pt] ($(RM-1-3.north east)+(0.2,-0.1)$) -- ++(0.2,0) |- ($(RM-3-3.south east)+(0.2,0.1)$);

% 4. The Path (The Zigzag Line)
% Draw connecting lines between cell anchors
\begin{scope}[rm_color, line width=1pt, line cap=round]
    % Row 1
    \draw (RM-1-1) -- (RM-1-2) -- (RM-1-3);
    % Row 2 Jump and Sweep
    \draw (RM-1-3) -- (RM-2-1);
    \draw (RM-2-1) -- (RM-2-2) -- (RM-2-3);
    % Row 3 Jump and Sweep
    \draw (RM-2-3) -- (RM-3-1);
    \draw (RM-3-1) -- (RM-3-2);
    % Final arrow segment, offset slightly from the last element a33
    \draw (RM-3-2) -- (RM-3-3);
\end{scope}

% ==========================================
% RIGHT: COLUMN-MAJOR TRAVERSAL
% ==========================================
% 1. The Matrix Grid (Placed right of the first matrix)
\matrix (CM) [matgrid, right=2.0cm of RM] {
    a_{00} & a_{01} & a_{02} \\
    a_{10} & a_{11} & a_{12} \\
    a_{20} & a_{21} & a_{22} \\
};

% 2. The Title (Placed relative to the matrix)
\node[title, above=0.3cm of CM-1-2] (cm_title) {Column-major order};

% 3. The Math Brackets (drawn manually)
\draw[line width=1.5pt] ($(CM-1-1.north west)+(-0.2,-0.1)$) -- ++(-0.2,0) |- ($(CM-3-1.south west)+(-0.2,0.1)$);
\draw[line width=1.5pt] ($(CM-1-3.north east)+(0.2,-0.1)$) -- ++(0.2,0) |- ($(CM-3-3.south east)+(0.2,0.1)$);

% 4. The Path (The N-shape Lines)
\begin{scope}[cm_color, line width=1pt, line cap=round]
    % Col 1
    \draw (CM-1-1) -- (CM-2-1) -- (CM-3-1);
    % Col 2 Jump and Sweep
    \draw (CM-3-1) -- (CM-1-2);
    \draw (CM-1-2) -- (CM-2-2) -- (CM-3-2);
    % Col 3 Jump and Sweep
    \draw (CM-3-2) -- (CM-1-3);
    \draw (CM-1-3) -- (CM-2-3);
    % Final arrow segment, offset from a33
    \draw (CM-2-3) -- (CM-3-3);
\end{scope}

\end{tikzpicture}

\end{figure}

This implementation chooses row-major as the data representation for our matrices. The index of the element $a_{i,j}$ of
an $n\times m$ matrix $A$ in the row-major array representing the matrix is $im + j$.

We can implement a basic matrix multiplication algorithm based on the dot-product formulation from Definition \ref{def:mat-mul}.

\begin{lstlisting}[mathescape=true]
function multiply(Matrix<n,m> X, Matrix<m,p> Y) {
  Matrix<n,p> V;

  for (i : [0,...,n-1]) {
    for (j : [0,...,p-1]) {
      int sum = $-\infty$;
      for (k : [0,...,m-1]) {
        sum = max(sum, X[i,k] + Y[k,j]);
      }
      V[i,j] = sum;
    }
  }

  return V;
}
\end{lstlisting}

Modern CPU caching exploits spatial locality of reference, where if a piece of data is accessed by the CPU, other pieces
of data in close proximity will be in the same cache line or prefetched into the cache, therefore reducing the chance of
an expensive cache miss.
In the chosen algorithm, matrices X and Y are both stored in row-major order, yet inside the inner loop, X is being accessed in
row-major order, and Y is being accessed in column-major order. The elements of Y being accessed on consecutive iterations 
are not in close proximity to each other in memory, and so cache misses dominate the running time, rather than the 
actual calculations.

This issue can be avoided by storing matrices in both row-major and column-major order. Then, depending on whether the
matrix is the right or left operand in a multiplication, we can select the most cache-friendly option to traverse.
This however, doubles our memory usage, and so we instead try to change the order in which we access elements of the matrices.

\subsection{Traversal order and cache performance}

The three loops in the previous matrix multiplication can be arbitrarily swapped with each other without an effect 
on correctness or asymptotic running time. We can swap our $k$ and $j$ loops so that the inner loop traverses Y in 
row-major order, resulting in the following algorithm:

\begin{lstlisting}[mathescape=true]
function multiply(Matrix<n,m> X, Matrix<m,p> Y) {
  Matrix<n,p> V;

  for (i : [0,...,n-1]) {
    for (k : [0,...,m-1]) {
      int xVal = X[i,k];
      if (xVal == $-\infty$) continue;

      for (j : [0,...,p-1]) {
        V[i,j] = max(V[i,j], xVal + Y[k,j]);
      }
    }
  }

  return V;
}
\end{lstlisting}

Apart from the improved cache performance, there are three other optimisations.

\begin{description}
  \item[Value hoisting] Since the indices $i$ and $k$ do not change in the inner loop of the algorithm, we can hoist the 
    retrieval of $X_{i,k}$ outside of the loop.
  \item[Sparsity] This optimisation concerns line 7 of the algorithm. To make sense of it, we can consider an extension
    to the previous matrix multiplication definition:
    $$
    v_{i \cdot} = \displaystyle\sum_{k=1}^{n} \! {}_{\oplus} x_{ik} \odot y_{k\cdot}
    $$
    This says that the $i$th row of $V$ is a linear combination of the $k$ rows of $Y$. More importantly, note that 
    $$
    -\infty \odot
    \begin{bmatrix}
      y_{k 1} \\
      \vdots \\
      y_{k p}
    \end{bmatrix}
    =
    \begin{bmatrix}
      -\infty \\ 
      \vdots \\ 
      -\infty
    \end{bmatrix}
    $$
    and since $-\infty$ is the additive identity, this row does not contribute to the final sum, and so the dot product 
    can be skipped. So, more precisely, the $i$th row of $V$ is a linear combination of the $k$ rows of $Y$ for which 
    $x_{ik}$ is not $-\infty$ \cite{gustavson-sparse}. The matrices in this project tend to be fairly sparse and so this 
    results in a significant performance improvement.
  \item[Vectorisation] Finally, since $Y$ is traversed contiguously in the inner loop, this makes it easily open to 
    optimisations using vector instructions. Vector instructions, also known as SIMD (Single-Instruction, Multiple-Data) \cite{flynntaxonomy},
    perform the same operation on several operands, exploiting data-level parallelism. The max and addition operations can 
    then be performed on the adjacent data from several iterations of the inner loop using vectorised variants, increasing 
    the number of elements processed per iteration.
\end{description}

\subsection{Matrix tiling}

%TODO: Do i need a more in-depth explanation?

One other common method of improving the cache performance of matrix multiplication algorithms is \textit{tiling}.
This involves implicitly dividing the matrix into square tiles of smaller sizes, and instead of traversing an entire 
row and column of the operands in one sweep to compute a cell of the output matrix, the algorithm works only on the values 
within one tile, before moving onto the next. This improves the spatial locality of memory accesses, as elements at the 
beginning of a row are not evicted by the time we get to the end of the row, which is a problem for large matrices.

However, since we are only tracking register dependencies in this project, our matrices are only approximately $64 \times 64$
in size. Assuming each cell holds a 32-bit integer, an entire matrix is stored in $64\cdot 64 \cdot 4 =16384$ bytes, 
which can be easily held in the 32 or 64KiB L1 cache of a modern processor core.

In this case, tiling the matrix will actually result in unnecessary overhead from a more complex control flow.

\subsection{Sparse Matrices}



%TODO: Diagram


\subsubsection{CSR Matrix Format}




\subsubsection{Sparse Matrix Multiplication}

\section{Expression Reduction}

\subsection{Fast exponentiation}

Can do fast exponentiation logarithmically. However repetitions usually happen on the loop level, not on the instruction level, so we need to first shrink the representation of a loop into a single matrix.

\subsection{Pair-reduction}

An approach that ends up reducing loops into a single matrix is a general one - memoisation.

%TODO: #check fixed point
This leads to an iterative (FIXED POINT? ) computation, where at each iteration, we select a pair of matrices to compute their product. We then cache that value, and scan through the entire expression, replacing every instance of that pair with the computed result. After that pass, we look for any exponents in the expression to reduce logarithmically.

At each iteration, we will shrink a pair of matrices in all iterations of a loop's body into a single matrix. So, after several iterations, the loop's body will be shrunk into a singular matrix, and then we can leverage logarithmic exponentiation to quickly shrink the repeated loop matrices into a single matrix.
%TODO: #todo include diagram demonstrating

We want to make sure that we make the most cost-effective multiplication at each iteration. Clearly we save the most cost by making the most frequent calculation in the matrix. 
To ensure this, we iterate through the expression, and select the most frequent pair of them all.
%TODO: #clarify on why we absolutely need this! see below para.

This also has the side-effect of reducing common pairs of instructions quickly through memoisation of the result. 

Once the frequency of the most frequent pair is 1, we can safely continue to directly multiplying matrices (as now there is no point in memoising values or doing frequency computations).

%TODO: todo add pseudocode

\subsection{Multiple pair-reduction}

*while the single-pair reduction ensures the minimum number of multiplications are made, frequency calculation is expensive. As the algorithm progresses, the frequency of pairs tends to reduce, and so the cost of the frequency calculation rises in comparison to the cost saved by memoising the most frequent pair.
We can make a tradeoff, replacing more pairs per iteration, which means we make less iterations and less frequency calculations*  %TODO: #rewrite 

*doesn't respect loop structure/why it doesn't ensure minimum number of multiplications*

Consider cksum's inner loop of instructions. They form a 7-instruction cycle, repeating 65536 times, that looks something like this:
$$
(A B C D E F G)^{65536}
$$
Now run the expression reduction, replacing the top 8 most frequent pairs:
%TODO: #todo add drawn diagram
We can see that we have reduced two iterations of the loop into one cycle of 7 matrices repeating 32768 times.
Since this is not exponentiation of one singular matrix, we cannot reduce this using logarithmic exponentiation, and instead have to go through more iterations of frequency computation and pair reductions (at each iteration halving the exponent) until we end up with a single loop of 7 matrices with an exponent of 1.
%TODO: #justify do the maths on which is better ^^^ #rewrite 
We haven't respected the boundaries of the loop and therefore haven't reduced INSIDE the loop

%TODO: #todo variable N-pair processing!!!

We want to exploit both the optimal reduction of replacing fewer pairs per iteration, and the increased throughput of replacing more pairs.
Because of our frequency-based approach, loop bodies shrink earlier (as the body's instructions occur many times!). Once the loops are reduced, we then have lower frequency values which we want to deal with by reducing more pairs every iteration.
We can do this by increasing the number of pairs to replace each iteration. How to do this, I am unsure as I haven't actually implemented any code. How silly of me.

%TODO: #todo analysis of how equation size changes with paircount

\subsection{Block preprocessing}

In the data given to us we can actually see branch structure so we reduce with respect to loops better - some loops occur over many different branches - this doesn't help here.

\subsection{Finding loops}

\section{Thread-level parallelism}

\section{Accelerator Compute}

